\subsection{The topology of compact convergence on the dual of a Fréchet space}

THEOREM 1 (Banach-Dieudonné). --- Let E be a locally convex metrizable space. The following topologies coincide on the dual E' of E :

\begin{enumerate}
\def\labelenumi{\alph{enumi})}
\item
  the topology \(\mathcal{T}_{\mathfrak{N}}\) of \(\mathfrak{N}\) -convergence, where \(\mathfrak{N}\) is the family of subsets of \(E\) each consisting of points of a sequence converging to 0;
\item
  the topology \(\mathcal{T}_{\mathrm{c}}\) of uniform convergence on compact subsets of E;
\item
  the topology \(\mathcal{T}_{pc}\) of uniform convergence on precompact subsets of E;
\item
  the topology \(\mathcal{T}_f\) which is the finest topology inducing the same topology as \(\sigma(\mathrm{E}',\mathrm{E})\) on every equicontinuous subset of \(\mathrm{E}'\) .
\end{enumerate}

First observe that a subset A of \(E'\) is closed for \(T_{f}\) if and only if \(A \cap H\) is closed for \(\sigma(E', E)\) for every subset H of \(E'\) which is equicontinuous and closed for \(\sigma(E', E)\) . The weak topology \(\sigma(E', E)\) and \(T_{pc}\) induce the same topology on every equicontinuous subset of \(E'\) (III, p.~17, prop. 5). Consequently each of the topologies \(T_{R}\) , \(T_{c}\) , \(T_{pc}\) , \(T_{f}\) is coarser than the one following it. It is therefore enough to prove that \(T_{R}\) is finer than \(T_{f}\) . Moreover, every translation in \(E'\) is a homeomorphism for \(T_{f}\) . Hence it is enough to prove that, if F is a subset of \(E'\) which is closed for \(T_{f}\) , and does not contain 0, then there exists a set \(S \in R\) such that \(S^{\circ} \cap F = \varnothing\) .

Let \((\mathbf{U}_{n})_{n\geqslant0}\) be a decreasing sequence of neighbourhoods of 0 in E forming a fundamental system of neighbourhoods of 0. We shall construct, by induction on \(n\geqslant0\) , finite sets \(X_{n}\) such that we have

\[
\mathrm{X} _ {n} \subset \mathrm{U} _ {n}\tag{4}
\]

\[
\left(\bigcup_ {0 \leqslant p \leqslant n} \mathrm{X} _ {p}\right) ^ {\circ} \cap \mathrm{U} _ {n + 1} ^ {\circ} \cap \mathrm{F} = \varnothing\tag{5}
\]

for every integer \(n \geqslant 0\) . Let \(m \geqslant 0\) be an integer such that \(X_{n}\) has been constructed for \(0 \leqslant n < m\) and satisfies (4) and (5) for \(0 \leqslant n < m\) . For every \(x \in U_{m}\) , put

\[
\mathrm{F} _ {x} = \big (\bigcup_ {0 \leqslant p <   m} \mathrm{X} _ {p} \big) ^ {\circ} \cap \{x \} ^ {\circ} \cap \mathrm{U} _ {m + 1} ^ {\circ} \cap \mathrm{F}.
\]

Formula (5) with \(n = m - 1\) implies that \(\bigcap_{x\in \mathrm{U}_m}\mathrm{F}_x = \emptyset\) . Further, the set \(\mathrm{U}_{m + 1}^{\circ}\) is equicontinuous, and compact for \(\sigma (\mathrm{E}',\mathrm{E})\) . In view of the definition of \(\mathcal{T}_f\) , each of the sets \(\mathrm{F}_x\) is compact for \(\sigma (\mathrm{E}',\mathrm{E})\) ; therefore there exists a finite subset \(\mathrm{X}_m\) of \(\mathrm{U}_m\) such that \(\bigcap_{x\in \mathrm{U}_m}\mathrm{F}_x = \emptyset\) , i.e.~relation (5) is satisfied for \(n = m\) .

Put \(S = \bigcup_{n\geqslant 0}X_n\) . We have \(X_{n}\subset U_{p}\) for \(n\geqslant p\) , therefore \(S\) is the set of points of a sequence which converges to 0 in \(E\) . From (5) we deduce that \(S^{\circ}\cap U_{n + 1}^{\circ}\cap F = \emptyset\) , and since \(E'\) is the union of the sequence of sets \(U_{n + 1}^{\circ}\) , we get \(S^{\circ}\cap F = \emptyset\) .

COROLLARY 1. --- Let E be a locally convex metrizable space. Every precompact subset of E is contained in the closed convex balanced envelope of the set of points of a sequence converging to 0.

This follows from the fact that the topologies \(\mathcal{T}_{pc}\) and \(\mathcal{T}_{\mathfrak{N}}\) are identical, on account of prop. 2 of III, p.~15.

COROLLARY 2. --- Let E be a Fréchet space. In order that a convex subset A of the dual E' of E be closed for \(\sigma(E', E)\) , it is necessary and sufficient that \(A \cap U^{\circ}\) is closed for \(\sigma(E', E)\) for every neighbourhood U of 0 in E.

Since E is complete, the topology \(T_{c}\) on \(E'\) is compatible with the duality between \(E'\) and E (IV, p.~3, Example); consequently the closed convex subsets in \(E'\) are the same for \(T_{c}\) and \(\sigma(E', E)\) (IV, p.~1, prop. 1). The corollary then follows from the identity of the topologies \(T_{c}\) and \(T_{f}\) .

Recall (I, p.~13) that the hyperplanes of E' which are closed for \(\sigma(\mathrm{E}', \mathrm{E})\) are the kernels of linear forms on E' associated with elements of E. Cor. 2 therefore gives another proof (for Fréchet spaces) of cor. 1 of III, p.~21.

COROLLARY 3. --- Let E be a Banach space and M a vector subspace of the dual E' of E. In order that M be closed for the weak topology \(\sigma(E', E)\) , it is necessary and sufficient that its intersection with the unit ball (closed) in E' be closed for \(\sigma(E', E)\) .

Example. --- * Let H be a hilbertian space satisfying the first axiom of countability; let \(H_{\sigma}\) denote the space H with the weakened topology assigned to it. Let \(\mathcal{L}^{1}(H)\) be the Banach space of nuclear endomorphisms of H (V, p.~51, and TS, V); the norm in \(\mathcal{L}^{1}(H)\) is defined by \(\|u\|_{1} = \operatorname{Tr}((u^{*}u)^{1/2})\) . We can identify \(\mathcal{L}(H)\) with the dual of the Banach space \(\mathcal{L}^{1}(H)\) by associating the linear form \(\phi_{u}: v \mapsto \operatorname{Tr}(uv)\) on \(\mathcal{L}^{1}(H)\) with every \(u \in \mathcal{L}(H)\) . Let A be a sub-algebra of \(\mathcal{L}(H)\) , containing 1 and stable under \(u \mapsto u^{*}\) ; this is a von Neumann algebra if and only if it is closed in \(\mathcal{L}(H)\) for the weak topology \(\sigma(\mathcal{L}(H), \mathcal{L}^{1}(H))\) . From cor. 3, we deduce the following criterion: for A to be a von Neumann algebra, it is necessary and sufficient that if \((u_{n})\) is any sequence of elements of A with norm \(\leqslant 1\) having a limit u in the space \(\mathcal{L}_{s}(H; H_{\sigma})\) , then u belongs to A.

